\section*{历史注记}
\phantomsection
\addcontentsline{toc}{section}{历史注记}

（方括号中的数字指本注记末尾的参考文献。）

在十九世纪初，任意函数的概念几乎无人知晓。更不用说，一般地研究函数的集合并给它们赋予拓扑结构的思想，在黎曼的时代之前并未出现（见第 I 章的历史注记），而且直到十九世纪末，这一思想才得到系统的运用。

然而，实值函数序列收敛的思想，自从无穷小分析诞生以来就或多或少地被自觉不自觉地使用着。当然，这里所谓收敛指的是逐点收敛；其他的收敛类型，要等到收敛级数与连续函数的概念由波尔查诺和柯西精确加以定义之后，才可能被描述。后者起初并未认识到逐点收敛与一致收敛之间的区别，并且相信自己已经证明，任何连续函数的收敛级数之和都是连续的（{[}1{]}, (2), vol.~3, p.~120）。这个错误几乎立即就被阿贝尔指出，他同时用一个经典的论证表明，每个幂级数在其收敛区间内部都是连续的；这个论证在上述特殊情形中本质上使用了一致收敛的概念（{[}2{]}, pp.~223-224）。剩下的事只是把这一概念一般地表述出来，这项工作由斯托克斯和赛德尔于 1847-1848 年独立完成，柯西本人也于 1853 年完成（{[}1{]}, (1), vol.~12, p.~30）(*)。

在魏尔斯特拉斯与黎曼的影响下，一致收敛概念及有关问题的系统研究，在十九世纪最后三分之一的时间里由德国学派（汉克尔、杜·布瓦–雷蒙）、而尤其是由意大利学派发展起来；狄尼与阿尔泽拉明确了连续函数序列的极限为连续的必要条件，而阿斯科利则引入了等度连续这一基本概念，并证明了刻画连续函数紧集的定理{[}3{]}（这个定理后来由蒙泰尔在其“正规族”理论中加以普及，正规族即解析函数的相对紧集）。

另一方面，魏尔斯特拉斯本人发现了 \([4b]\) 在有界集上用多项式一致逼近一个或多个实变量的连续实值函数的可能性。这一结果立即引起了浓厚的兴趣，并引出了许多“定量的”研究(*)。

对这些问题，现代的贡献主要在于赋予它们以所能达到的完全一般性：考虑定义域与值域不再限于 R 或有限维空间的函数，从而借助一般拓扑概念把它们安置在其固有的环境中。特别地，魏尔斯特拉斯定理在经典分析中早已显示出是一件强有力的武器，近年来 M. H. 斯通把它推广到了广泛得多的场合；他发展了 H. 勒贝格引进的一个思想（在魏尔斯特拉斯定理的一个证明中），清楚地指明了函数格在连续实值函数逼近理论中所起的重要作用（用“格多项式”逼近，参看 § 4 的命题 2 与定理 2），并且表明了推广的魏尔斯特拉斯定理如何以一整系列类似的逼近定理作为其直接推论，这些定理因而能够以连贯得多的方式归拢在一起。我们或多或少遵循了他的阐述{[}5{]}。

